\documentclass[tikz,border=12pt]{standalone}
\usepackage{amsmath}
\usepackage{kotex}
\usetikzlibrary{arrows.meta}
\definecolor{surface}{RGB}{35,100,175}
\definecolor{piplane}{RGB}{135,75,160}
\definecolor{planestress}{RGB}{205,105,25}
\definecolor{hydro}{RGB}{30,130,95}
\begin{document}
\begin{tikzpicture}[>=Latex,font=\sffamily,line cap=round,line join=round]
% +h 방향에서 원점을 본다. h=(1,1,1)/sqrt(3)은 화면 밖을 향한다.
% 화면의 직교정규 기저는 u=(-1,1,0)/sqrt(2), v=(-1,-1,2)/sqrt(6)이다.
% 화면 좌표=(u:sigma, v:sigma)/sigmaY이며 배율은 4 cm이다.
% 투영된 세 단위축의 방향은 210도, 330도, 90도이고 길이는 모두 sqrt(2/3)이다.
\pgfmathsetmacro{\rad}{4*sqrt(2/3)}
\node[font=\large\bfseries] at (2.5,6.0) {정수압 축 방향에서 본 von Mises 항복면};
\fill[piplane!6] (-5,-4.95) rectangle (5.05,4.95);
\node[piplane,anchor=north west] at (-4.8,4.75) {$\pi$-plane: 화면과 나란한 평면};
\fill[surface!7] (0,0) circle (\rad);
\draw[black!65,->,thick] (0,0)--({4.4*cos(210)},{4.4*sin(210)})
 node[below left,align=right] {$\sigma_1$축의\\투영 방향};
\draw[black!65,->,thick] (0,0)--({4.4*cos(330)},{4.4*sin(330)})
 node[below right,align=left] {$\sigma_2$축의\\투영 방향};
\draw[black!65,->,thick] (0,0)--(0,3.9)
 node[above right,align=left] {$\sigma_3$축의\\투영 방향};
\draw[piplane,line width=3pt] (0,0) circle (\rad);
\draw[planestress,line width=1.1pt,dash pattern=on 5pt off 4pt] (0,0) circle (\rad);
\draw[black!55,->] ({1.05*cos(210)},{1.05*sin(210)}) arc[start angle=210,end angle=330,radius=1.05];
\node at (0,-1.43) {$120^\circ$};
\fill[planestress] ({-4/sqrt(6)},{-4/sqrt(2)}) circle (2.8pt);
\node[planestress,below left=4pt] at ({-4/sqrt(6)},{-4/sqrt(2)}) {$B$의 투영};
\fill[planestress] ({4/sqrt(2)},{-4/sqrt(6)}) circle (2.8pt);
\node[planestress,above right=5pt] at ({4/sqrt(2)},{-4/sqrt(6)}) {$Q$의 투영};
\draw[hydro,very thick,fill=white] (0,0) circle (.16);
\fill[hydro] (0,0) circle (2pt);
\node[below left=6pt] at (0,0) {$O$};
\node[surface,align=center] at (-1.3,1.4) {항복 원통의 내부\\$f<0$};
\node[piplane,anchor=west,align=left] at (6.0,3.5) {보라색 원: 원통과 $\pi$-plane의 교선\\실제 3차원 공간에서도 원이다.};
\node[planestress,anchor=west,align=left] at (6.0,1.7) {주황색 점선: $\sigma_3=0$ 교선의 투영\\앞의 주황색 타원을 정수압 축 방향으로\\투영하면 보라색 원과 겹친다.};
\node[hydro,anchor=west,align=left] at (6.0,-.35) {녹색 $\odot$: 정수압 축의 양의 방향\\$\boldsymbol h=(1,1,1)/\sqrt3$은 화면 밖을 향한다.};
\node[anchor=west,align=left] at (6.0,-2.35) {세 축의 투영 방향 사이 각도는 모두 $120^\circ$다.\\이는 원래 주응력 축 사이의 $90^\circ$와 다르다.\\원래의 세 축은 서로 수직이다.};
\node[anchor=west,align=left] at (6.0,-4.15) {원 반지름: $r=\sqrt{2/3}\,\sigma^Y$\\원점으로부터의 거리는 편차응력의 크기다.};
\node[align=center,font=\small,text=black!70] at (2.8,-5.8) {응력점에서 평균응력 성분을 빼면 $\pi$-plane으로 투영된다.\\$\boldsymbol s=\boldsymbol\sigma-\sigma_m\boldsymbol I$: 정수압 축에 평행한 이동};
\end{tikzpicture}
\end{document}
